\section{Week 11}


\subsection{Opening Remarks with Exponential of a Matrix}

\begin{remark}
    For which matrices is it easy to compute the exponential?
    \begin{enumerate}
        \item[(1)] Diagonal matrices
        \item[(2)] Nipoltent matrices (finite sum); that is, matrices where there exists a \( k \in \N  \) such that \( A^{k} = O_{n \times n } \). 
            \begin{itemize}
                \item An example,  
                    \[  A = \begin{pmatrix} 1 & 1 \\ - 1 & - 1  \end{pmatrix}  \]
                    is nilpotent.
                \item More generally, nilpotent \( 3 \times 3  \) matrices have the following form
                    \[  \begin{pmatrix} a & a & a \\ b & b & b \\  (a+b) & -(a+b) & -(a+b) \end{pmatrix}  \]
                    is nilpotent.
            \end{itemize}
        \item[(3)] Diagonalizable matrices
    \end{enumerate}
\end{remark}

\begin{prop}[Powers and Exponential of Diagonalizable Matrices]
   Suppose that \( A \in M(n,\C)  \) is \textbf{Diagonalizable}, that is, suppose there exists an invertible matrix \( P  \) and a diagonal matrix \( D \) such that  
   \[  A = P D P^{-1}. \]
   Then 
   \begin{enumerate}
       \item[(i)] For all \( m \in \N  \) \( A^{m} = P D^{m} P^{-1} \).
        \item[(ii)] \( \text{exp}(A) = P (\text{exp}(D)) P^{-1} \) 
   \end{enumerate}
\end{prop}
\begin{proof}
\begin{enumerate}
    \item[(i)] For \(  m = 1  \), \( A^{1} = P D^{1} P^{-1}  \). For \( m = 2  \), \[ A^{2} = AA = P D P^{-1} (P D P^{-1})  = P D D P^{-1} = P D^{2} P^{-1}. \] We can do a similar argument by using induction.
        \item[(ii)] By definition, we have
            \begin{align*}
                e^{A} = \sum_{  m = 0  }^{ \infty  } \frac{ A^{m} }{ m!  }  &= \sum_{ m = 0  }^{ \infty   } \frac{ (P D P^{-1})^{m} }{ m !   }  \\
                      &= \sum_{ m = 0  }^{ \infty   } \frac{ P D^{m} P^{-1} }{ m!  }  \\
                      &= P \Big(  \sum_{ m = 0  }^{ \infty  } \frac{ D^{m} }{ m! }  \Big) P^{-1}  \\ 
                      &= P e^{D} P^{-1}.
            \end{align*}
\end{enumerate}
\end{proof}

\begin{remark}[Review of Linear Algebra]
   If the matrix \( A \in M(n,\C) \) has \( n  \) linearly independent eigenvectors, then \( A  \) is diagonalizable. Indeed, 
   \[  A = P D P^{-1} \]
   where 
   \begin{enumerate}
       \item[(*)] Columns of \( P \) are linearly independent eigenvectors of \( A  \)
        \item[(*)] Diagonal entries of \( D  \) are the corresponding eigenvalues.
   \end{enumerate}
\end{remark}

\subsection{Review of Eigenvectors}

\begin{eg}
    If \( A = \begin{pmatrix} 3 & 1 \\ 1 & 3  \end{pmatrix}  \), the eigenvalues are \( 4  \) and \( 2  \) which correspond to the (linearly independent) vectors \( (1,1) \) and \( (-1,1) \). Since the number of linearly independent vectors is \( 2 \). Then 
    \begin{align*}
        A = P D P^{-1} \  &\text{where} \ P = \begin{pmatrix} 1 & - 1 \\ 1 & 1  \end{pmatrix} \ \text{and} \ D = \begin{pmatrix} 4 & 0 \\ & 0 & 2  \end{pmatrix}   \\
    \end{align*}
    where 
    \[  A \begin{pmatrix}  1 \\ 1  \end{pmatrix}  = 4 \begin{pmatrix}  1 \\ 1  \end{pmatrix}  \] and
    \[  A \begin{pmatrix} - 1 \\ 1  \end{pmatrix}  = 2 \begin{pmatrix}  - 1 \\ 1  \end{pmatrix}. \]
    Also, it is true tat 
    \[  A = P D P^{-1} \]
    where
    \[  P = \begin{pmatrix}  - 1 & 1 \\ 1 & 1  \end{pmatrix} , D = \begin{pmatrix} 2 & 0 \\ 0 & 4  \end{pmatrix}. \]
\end{eg}

\begin{eg}
   Let  
   \[  A = \begin{pmatrix} 2 & 1 & 1 \\ 2 & 3 & 2 \\ 3 & 3 & 4  \end{pmatrix}  \] 
   where the eigenvalues and the corresponding eigenvectors are \( 1  \) and \(  7  \) with 
   \( (-1,1,0) \) and \( (-1,0,1) \) (corresponding to \( 1 \)) and \( 7  \) corresponding to \( (1,2,3) \). Then since there are \( 3  \) linearly independent eigenvectors, we have
   \[  P = \begin{pmatrix} -1 & -1 & 1 \\ 1 & 0 & 2 \\ 0 & 1 & 3  \end{pmatrix}  \]
   and
   \[  D = \begin{pmatrix} 1 & 0 & 0 \\ 0 & 1 & 0 \\ 0 & 0 & 7 \end{pmatrix}.  \]
\end{eg}

In summary, in order to compute \( e^{X} \) where \( X  \) is diagonalizable requires the following steps:
\begin{enumerate}
    \item[(1)] Write \( X = P D P^{-1} \)
    \item[(2)] Then \( e^{X} = P e^{D} P^{-1} \).
\end{enumerate}

\subsection{Computing Exponential of Matrices}

\begin{eg}[Skew Symmetric Matrices and its Exponential]
    Let \( \theta \in \R  \) be a fixed number. Let \( X = \begin{pmatrix} 0 & - \theta \\ \theta & 0  \end{pmatrix}  \). Prove that  
    \[  e^{X} = \begin{pmatrix}  \cos \theta & - \sin \theta \\ \sin \theta & \cos \theta \end{pmatrix}. \]
\end{eg}
\begin{proof}
The corresponding eigenvalues and eigenvectors of this matrix are \( - i \theta  \) corresponding to \( (1,i) \) and \( i \theta  \) corresponding to \( (i,1) \). Check this in private! So, \( X = P D P^{-1} \) where 
\[  P = \begin{pmatrix} 1 & i \\ i & 1  \end{pmatrix}  , \ D = \begin{pmatrix} -i \theta & 0 \\ 0 & i \theta \end{pmatrix}. \]
Note that 
\[  P^{-1} = \frac{ 1 }{ 1 - i^{1} }  \begin{pmatrix} 1 & - i \\ - i & 1  \end{pmatrix}  = \frac{ 1 }{ 2 }  \begin{pmatrix} 1 & - i \\ - i & 1  \end{pmatrix} = \begin{pmatrix} \frac{ 1 }{ 2 }  & \frac{ -i }{ 2 } \\ \frac{ -i }{ 2 }  & \frac{ 1 }{ 2 }  \end{pmatrix}.  \]
Therefore,
\[  e^{X} = P e^{D}p^{-1} = \begin{pmatrix}  1 & i \\ i & 1  \end{pmatrix} \begin{pmatrix}  e^{-i \theta} & 0 \\ 0 & e^{i \theta} \end{pmatrix} \begin{pmatrix} \frac{ 1 }{ 2 }  & \frac{ -i }{ 2 }   \\ \frac{ -1 }{ 2 }  & \frac{ 1 }{ 2 } \end{pmatrix}.  \]
Then we have
\begin{align*}
    e^{X} &= \begin{pmatrix} e^{-i \theta} & i e^{i \theta} \\ i e^{-i \theta} & e^{i \theta} \end{pmatrix}  \begin{pmatrix}  \frac{ 1 }{ 2 }  & \frac{ -i  }{ 2 }  \\ \frac{ -i  }{ 2 } & \frac{ 1 }{ 2 }  \end{pmatrix} \\
          &= \begin{pmatrix}  \frac{ e^{- i \theta} + e^{i \theta} }{  2  }  & \frac{ -i e^{=i \theta} + i e^{i \theta} }{ 2 } \\ \frac{ i e^{-i \theta} - i e^{i \theta} }{ 2  } & \frac{ e^{-i \theta} + e^{i \theta} }{ 2  }   \end{pmatrix}  \\
          &= \begin{pmatrix} \cos \theta & - \sin \theta \\ \sin \theta & \cos \theta \end{pmatrix}
\end{align*}
where 
\[ \cos \theta = \frac{ e^{i \theta} + e^{- i \theta}  }{  2  }  \ \text{and} \ \sin \theta =  \frac{ e^{i \theta} - e^{-i \theta} }{ 2 i }. \]
Therefore, we have if we have \( X = \begin{pmatrix}  0 & - \theta \\ \theta & 0  \end{pmatrix}  \), then
\[  e^{tx} = \begin{pmatrix} \cos (\theta t) & - \sin (\theta t) \\ \sin (\theta t) & \cos (\theta t) \end{pmatrix}  \ \forall t \in \R.  \]

\end{proof} 

\begin{remark}
    In general, every matrix \( X \in M(n,\C) \) can be written (uniquely) as
    \[  X = S  +N  \]
    where \( S  \) is diagonalizable, \( N  \) is nilpotent, and \( SN = NS  \). Then 
    \[  e^{X} = e^{S+N} = e^{S} e^{N} \]
    where \( e^{S} \) can be easily computed by diagonalization and \( e^{N} \) can be computed by converting the infinite sum to a finite sum.
\end{remark}
\begin{remark}
    Let \( G  \) be a Lie group and \( e \in G  \), be its identity element. Without loss of generality, we may assume that \( G  \) is an embedded submanifold of \( \R^{k} \) for some \( k \in \N \) (so, our construction of \( {T}_{e}G \) is valid). A question we can ask is "what sort of algebraic operations can we define on \( {T}_{e}G \)?"
\end{remark}
 
\subsection{Operations of Tangent Space at the identity of a Lie Group}

Let \( X,Y \in {T}_{e}G \). Then there exists a smooth path \( A(t) \) in \( G  \) such that \( A(0) = e  \) and \( A'(0) =  X \). Likewise, there exists another path \( B(t)  \) in \( G  \) such that \( B(0) = e  \) and \( B'(0) = Y  \). The first immediate operation we can do with these two paths is to multiply them together which in of itself is a smooth path because the operation is smooth. That is,  
\[  C(t) = A(t) B(t)  \]
is another smooth path in \( G  \) with \( C(0) = e  \) and so we have \( C'(0 ) \in {T}_{e}G \). In particular, if \( G  \) is a matrix Lie group, then
\[ C'(t) = A'(t) B(t) + A(t) B'(t) \implies C'(0) = A'(0) B(0) + A(0) B'(0) = X I + I Y = X  + Y.   \]
If \( X  \) and \( Y  \) are in \( {T}_{I}G  \), then \( X + Y  \in {T}_{I}G\). Hence, we now know that \( {T}_{e}G \) is closed under addition.

Now, we will show that this space is closed under scalar multiplication. For any \( \alpha \in \R  \), \( C(t) = A(\alpha t)  \) is another smooth path in \( G  \) with \( C(0) = e  \). Then \( C'(0) \in {T}_{e}G \). In particular, if \( G  \) is a matrix Lie group, then  
\begin{align*}
    C'(t) &= \alpha A'(\alpha t )  \\
          &= \alpha A'(0) \\
          &= \alpha X. 
\end{align*} 
So, we proved that if \( X \in {T}_{e}G \) and \( \alpha \in \R  \), then \( \alpha X \in {T}_{e} G  \). Hence, this tells us that \( {T}_{e}G  \) is a linear subspace of \( \R^{k} \). 

Note that \( C(t) = [A(t)]^{-1} \) is another smooth path in \( G  \) with \( C(0) = e  \). That is, \( C'(0) \in {T}_{e}G \). In particular, if \( G  \) is matrix Lie group. Then
\begin{align*}
    &C'(t) = - A^{-1}(t) A'(t) A^{-1}(t)   \\
    &\implies C'(0) = - A^{-1}(0) A'(0) A^{-1}(0) \\
    &= - I X I  \\
    &= - X.
\end{align*}
That is, if \( X \in {T}_{I} G  \), then \( - X \in {T}_{I}G \) (a special case of the last proof).

Next, if we are given \( E \in G  \) where \( C(t) = E B(t) \) is a smooth path in \( G  \) but \( c(0) \neq e  \). However, we may consider \( C(t)  = E B(t) E^{-1} \). Then \( C(0) = e  \). Indeed, 
\[  C(0) = E B(0) E^{-1} = E e E^{-1} = E E^{-1} = e.  \]
In particular, if \( G  \) is a matrix Lie group, then
\begin{align*}
    &C'(t) = E B'(t) E^{-1} \\
    &\implies C'(0) = E Y E^{-1}.
\end{align*}
That is, if \( Y \in {T}_{I} g  \), then for every matrix \( E \in G  \), \( E Y E^{-1} \in {T}_{I}G \).

For all \( s  \), \( {C}_{s}(t) = A(s) B(t) A(s)^{-1} \) is a smooth path in \( G  \) and for all \( s  \), \( {C}_{s}(0) = e  \). Then for all \( s  \), \( {C}_{s}'(0) \in {T}_{e} G  \). So, if we let \( v(s) = {C}_{s}'(0) \), then \( v(s) \) is a path inside \( {T}_{e}G  \). Thus, in particular, \( v'(0) \in {T}_{e}G \). That is, 
\[ \frac{ d  }{ ds  } \Big|_{s = 0} \frac{ \partial  }{  \partial t  } \Big|_{t =0} A(s) B(t) A^{-1}(s) \in {T}_{e}G.\]
Note that please notice that if \( G  \) is abelian, then the above expression is \( 0  \). Now, suppose \( G  \) is a matrix Lie group. Let's compute the above vector in \( {T}_{e}G \). We have
\begin{align*}
    \frac{ d }{ ds  }  \Big|_{s = 0} \frac{ \partial  }{ \partial t  }  \Big|_{ t = 0 } A(s) B(t) A^{-1}(s) &= \frac{ d }{ ds  }  \Big|_{s=0 } A(s) B'(0) A^{-1}(s)   \\
                                                                                                            &=  A'(0) B'(0) A^{-1}(0) + A(0) B'(0) (-A(0) A'(0) A^{-1}(0)) \\
                                                                                                            &= A'(0) B'(0) A^{-1}(0) - A(0) B'(0) ( A(0)A'(0) A^{-1}(0)) \\
                                                                                                            &= XYI - I Y(IXI) \\
                                                                                                            &= XY - YX.
\end{align*}
This remarkable equation that we have obtained at the end is called the \textbf{Lie Bracket}.

\begin{remark}[Observation]
    Suppose \( G  \) is an abstract group. Note that \( G  \) is abelian if and only if for all \( x,y \in G   \) \( xy = yx \). This is equivalent to for all \( x,y \in G  \), \( (xy)(yx)^{-1} = e  \) which is further equivalent to for all \( x,y \in G  \), \( xy (x^{-1} y^{-1}) = e  \).   
\end{remark}


Similarly, for all \( s  \), \( {C}_{s}(t) = A(s) B(t) A^{-1}(s) B^{-1}(t)  \) is a smooth path in \( G  \) and for all \( s  \), \( {C}_{s}(0) = e  \). Then we have for all \( s  \), \( {C}_{s}'  \in  {T}_{e}G\). Therefore, if we define \( v(s) = {C}_{s}'(0) \), then \( v(s) \) is a path inside the vector space \( {T}_{e}G \). Thus, \( v'(0) \in {T}_{e}G \). That is, 
\begin{align*}
    \frac{ d }{ ds  }  \Big|_{s = 0 } \frac{ d }{ dt  }  \Big|_{t = 0} A(s) B(t) A^{-1}(s) B^{-1}(t) \in {T}_{e}G \\
\end{align*}
Note that if \( G  \) is abelian, the above expression is zero. In particular, if \( G \) is a matrix Lie group, we have 
\begin{align*}
    \frac{ d }{ ds } \Big|_{s =0} \frac{ \partial  }{ \partial t  }  \Big|_{ t = 0 } [A(s) B(t) A^{-1}(s) B^{-1}(t) ] &= \frac{ d }{ ds } \Big|_{s = 0} [ A(s) B'(t) A^{-1}(s) B^{-1}(t) + A(s) B(t) A^{-1}(s) (- B^{-1}(t) B'(t) B^{-1}(t))  ] \Big|_{t = 0} \\
                                                                                                                      &= \frac{ d }{ ds }  \Big|_{s = 0} [A(s) B'(0) A^{-1}(s) B^{-1}(0) \\ 
                                                                                                                      &- A(s) B(0) A^{-1}(s) B^{-1}(0) B'(0) B^{-1}(0)] \\
                                                                                                                      &= \frac{ d }{ ds } \Big|_{s = 0} [ A(s) Y A^{-1}(s) - A(s) A^{-1}(s) Y ] \\
                                                                                                                      &= \frac{ d }{ ds } \Big|_{s = 0} [ A(s) Y A^{-1}(s) - Y ] \\
                                                                                                                      &= \frac{ d }{ ds } \Big|_{s = 0} [A(s) Y A^{-1}(s)] \\
                                                                                                                      &= XY - YX \tag{See previous proof}
\end{align*}

\begin{remark}
    Note that from the previous two cases produce the zero map on \( {T}_{e}G \) if \( G  \) is abelian:
    \[ \forall x ,y \in {T}_{e}G \ \ (x,y) \mapsto 0.  \]
\end{remark}

We define the \textbf{Lie Bracket} on \( {T}_{e}G \) by for all \( X,Y \in {T}_{e}G \)
\[  [X,Y] = \frac{ d }{ ds } \Big|_{s = 0} \frac{ \partial  }{ \partial t  }  \Big|_{t = 0} A(s) B(t) A^{-1}(s) B^{-1}(t).\]
When we equip \( {T}_{e}G  \) with this Lie bracket is called the Lie algebra of \( g  \) and is denoted by \( \text{Lie}(G) \).

\subsection{Lie Algebras}

\begin{definition}[Lie Algebra]
    A real (or complex) Lie algebra is a real (or complex) vector space \( L \) equipped with a function \( [\cdot, \cdot ] : L \times L \to L  \).
    \begin{enumerate}
        \item[(i)] \( [\cdot, \cdot]  \) is bilinear 
        \item[(ii)] \( [\cdot, \cdot]  \) is skew-symmetric
        \item[(iii)] \( [\cdot, \cdot]  \) satisfies the Jacobi identity; that is, for all \( x,y,z \in L \), we have 
            \[  [x,[y,z]] + [z,[x,y]] + [y, [z,x]] = 0.  \]
    \end{enumerate}
\end{definition}

One way to create a Lie Algebra is to start with a Lie group and look at the tangent space at the identity of the Lie group and define a Lie Bracket on it.

